\documentclass[12pt,a4paper]{report}

% ── Encoding & Fonts ──
\usepackage[utf8]{inputenc}
\usepackage[T1]{fontenc}
\usepackage{lmodern}

% ── Page Layout ──
\usepackage[margin=1in,headheight=15pt]{geometry}
\usepackage{setspace}
\onehalfspacing

% ── Language ──
\usepackage[english]{babel}

% ── Graphics & Colors ──
\usepackage{graphicx}
\usepackage[dvipsnames,svgnames,table]{xcolor}

% ── Math ──
\usepackage{amsmath,amssymb,amsfonts}

% ── Tables ──
\usepackage{booktabs}
\usepackage{tabularx}
\usepackage{longtable}
\usepackage{multirow}

% ── Lists ──
\usepackage{enumitem}

% ── Code Listings ──
\usepackage{listings}
\lstset{
  basicstyle=\ttfamily\small,
  keywordstyle=\color{blue}\bfseries,
  commentstyle=\color{gray},
  stringstyle=\color{red},
  breaklines=true,
  frame=single,
  backgroundcolor=\color{gray!5},
  numbers=left,
  numberstyle=\tiny\color{gray},
  numbersep=8pt,
  tabsize=4,
  showstringspaces=false,
}

% ── Hyperlinks ──
\usepackage[colorlinks=true,linkcolor=NavyBlue,citecolor=Green,urlcolor=RoyalBlue]{hyperref}

% ── Headers & Footers ──
\usepackage{fancyhdr}
\pagestyle{fancy}
\fancyhf{}
\fancyhead[L]{\leftmark}
\fancyhead[R]{\thepage}
\renewcommand{\headrulewidth}{0.4pt}

% ── Chapter Style ──
\usepackage{titlesec}
\titleformat{\chapter}[display]
  {\normalfont\huge\bfseries}
  {\chaptertitlename\ \thechapter}{20pt}{\Huge}
\titlespacing*{\chapter}{0pt}{-20pt}{30pt}

% ── Caption Styling ──
\usepackage{caption}
\captionsetup{font=small,labelfont=bf}

% ── Bibliography ──
\usepackage[numbers,sort&compress]{natbib}

% ── Misc ──
\usepackage{parskip}
\usepackage{float}

% ════════════════════════════════════════════════
%  DOCUMENT METADATA
% ════════════════════════════════════════════════
\title{\textbf{Parameterised CFD Simulation Guide}\\[6pt]
       \Large Hydrogen--Methane T-Junction Mixing\\
       in ANSYS Fluent with Automated Design-Point Sweeps}
\author{Yash}
\date{\today}

\begin{document}

% ── Title Page ──
\maketitle
\thispagestyle{empty}

% ── Table of Contents ──
\tableofcontents
\newpage


% ════════════════════════════════════════════════════════════════
%  CHAPTER 1 — Introduction and Problem Definition
% ════════════════════════════════════════════════════════════════
\chapter{Introduction and Problem Definition}
\label{ch:intro}

\section{Background}

Decarbonising domestic heating is a central pillar of net-zero strategies worldwide.
A promising transitional approach is to inject \textbf{green hydrogen} into the existing
natural-gas transmission network at volume fractions of 5--20\,\%, a range within
which most conventional boilers can operate without modification~\cite{Eames2022}.

The injection is typically performed through a \textbf{T-junction}---a branch pipe
that meets the main pipeline at an angle.  Because hydrogen is approximately
\textbf{eight times lighter} than methane (the dominant component of natural gas),
the injected jet lacks momentum to penetrate the main flow and is simultaneously
driven upward by buoyancy.  The result is \textit{stratification}: localised
hydrogen concentrations on the upper pipe wall that can exceed 40\,\% by volume,
far above the bulk target~\cite{Eames2022}.  Prolonged exposure of carbon-steel
pipe walls to such concentrations accelerates \textbf{hydrogen embrittlement},
reducing fracture toughness and enabling crack propagation~\cite{Laureys2022}.

Understanding and controlling the mixing process therefore requires
high-fidelity computational fluid dynamics (CFD) that resolves the transient
turbulent structures responsible for scalar transport.  This guide provides a
complete, literature-backed, step-by-step procedure to set up such a simulation
in \textbf{ANSYS Fluent}, with full parametric automation so that the influence
of key geometric and flow variables can be swept systematically.


\section{Parametric Variables}

Three independent parameters govern the T-junction mixing problem.
Their definitions, symbols, and proposed sweep ranges are summarised in
Table~\ref{tab:params}.

\begin{table}[H]
\centering
\caption{Parametric variables and their sweep ranges.}
\label{tab:params}
\begin{tabular}{@{}llll@{}}
\toprule
\textbf{Parameter} & \textbf{Symbol} & \textbf{Definition} & \textbf{Sweep Values} \\
\midrule
Diameter ratio    & $\beta$   & $d_2 / d_1$   & 0.25,\; 0.50,\; 0.75,\; 1.00 \\
Injection angle   & $\theta$  & Angle between branch axis and main axis & $30^\circ$,\; $60^\circ$,\; $90^\circ$,\; $120^\circ$,\; $150^\circ$ \\
Velocity ratio    & $R_v$     & $u_2 / u_1$   & 0.5,\; 1.0,\; 2.0,\; 4.0 \\
\bottomrule
\end{tabular}
\end{table}

\noindent To ensure robust parameterisation, the study is organised as
\textbf{five separate Workbench systems}---one per injection angle.
Within each system, only $\beta$ and $R_v$ are swept as Workbench
parameters.  This eliminates the topology changes that occur when the
injection angle varies (different branch--pipe intersection geometry),
guaranteeing that Named Selections, mesh settings, and boundary
conditions remain valid across all design points within each base.
Table~\ref{tab:bases} summarises the architecture.

\begin{table}[H]
\centering
\caption{Five-base architecture: one Workbench system per injection angle.}
\label{tab:bases}
\begin{tabular}{@{}llcl@{}}
\toprule
\textbf{System Name} & $\theta$ & \textbf{Design Points} & \textbf{Swept Parameters} \\
\midrule
\texttt{Theta\_030} & $30^\circ$  & 16 & $\beta \times R_v$ \\
\texttt{Theta\_060} & $60^\circ$  & 16 & $\beta \times R_v$ \\
\texttt{Theta\_090} & $90^\circ$  & 16 & $\beta \times R_v$ \\
\texttt{Theta\_120} & $120^\circ$ & 16 & $\beta \times R_v$ \\
\texttt{Theta\_150} & $150^\circ$ & 16 & $\beta \times R_v$ \\
\midrule
\multicolumn{2}{@{}l}{\textbf{Total}} & \textbf{80} & \\
\bottomrule
\end{tabular}
\end{table}

\noindent From these three independent parameters, several derived quantities
follow directly:

\begin{description}[style=nextline, font=\normalfont\bfseries, leftmargin=1.5em]

  \item[Volumetric dilution ($\mathcal{D}$):]
  The ratio of the volumetric flow rate through the branch to that through the
  main pipe~\cite{Eames2022}:
  \begin{equation}
    \mathcal{D} = \beta^{2}\, R_v
    \label{eq:dilution}
  \end{equation}

  \item[Momentum ratio ($M_R$):]
  The ratio of the dynamic pressure of the branch jet to that of the main
  stream, controlling jet penetration depth~\cite{Eames2022}:
  \begin{equation}
    M_R = \frac{\rho_2}{\rho_1}\, R_v^{2}
    \label{eq:momratio}
  \end{equation}
  Because $\rho_{\mathrm{H_2}} / \rho_{\mathrm{CH_4}} \approx 1/8$ at equal
  pressure and similar temperature, a velocity ratio of $R_v = 1$ yields
  $M_R \approx 0.125$---extremely weak penetration.

\end{description}


\section{Base Operating Conditions}

The simulation replicates the UK high-pressure gas transmission network
conditions used by Eames et al.~\cite{Eames2022}:

\begin{table}[H]
\centering
\caption{Base operating conditions (fixed across all design points).}
\label{tab:opcond}
\begin{tabular}{@{}lll@{}}
\toprule
\textbf{Quantity} & \textbf{Value} & \textbf{Source} \\
\midrule
Main-pipe diameter, $d_1$            & 460\,mm (18\,in)     & \cite{Eames2022} \\
Operating pressure, $p_{\text{op}}$  & 6.9\,MPa             & \cite{Eames2022} \\
Main-flow velocity, $u_1$            & 10\,m\,s$^{-1}$      & \cite{Eames2022} \\
Methane temperature, $T_1$           & 284\,K               & \cite{Eames2022} \\
Hydrogen temperature, $T_2$          & 293\,K               & \cite{Eames2022} \\
Main-flow species                    & Pure CH$_4$           & \cite{Eames2022} \\
Branch-flow species                  & Pure H$_2$            & \cite{Eames2022} \\
\bottomrule
\end{tabular}
\end{table}


\section{Simulation Strategy Overview}

The workflow proceeds through five major stages, each detailed in a dedicated
chapter.  The geometry and meshing stages (Chapters~\ref{ch:geom}
and~\ref{ch:mesh}) are repeated \textbf{once per injection angle} to create
the five base systems.  All subsequent stages (Fluent setup, solution,
parameterisation) are identical across bases and need only be configured once,
then duplicated.

\begin{enumerate}
  \item \textbf{Geometry} (Chapter~\ref{ch:geom}): Build the T-junction fluid
        domain in DesignModeler with a fixed injection angle and a
        parameterised branch diameter.  Repeat for each of the five angles.
  \item \textbf{Meshing} (Chapter~\ref{ch:mesh}): Generate a mesh with
        a parameterised Body of Influence and inflation layers targeting
        $y^+ \approx 1$.
  \item \textbf{Fluent Physics Setup} (Chapters~\ref{ch:models}--\ref{ch:bc}):
        Configure the solver, turbulence model (SA-DDES), species transport,
        material properties, and boundary conditions.
  \item \textbf{Solution and Execution} (Chapters~\ref{ch:solution}--\ref{ch:run}):
        Set numerical schemes, initialise via Hybrid Initialisation, and run the
        transient DDES directly from $t = 0$.
  \item \textbf{Parameterisation and Post-Processing}
        (Chapters~\ref{ch:param}--\ref{ch:post}): Sweep $\beta$ and $R_v$
        within each base system, then combine results across bases for
        analysis.
\end{enumerate}


% ════════════════════════════════════════════════════════════════
%  CHAPTER 2 — Parameterised Geometry
% ════════════════════════════════════════════════════════════════
\chapter{Parameterised Geometry (DesignModeler)}
\label{ch:geom}

\section{Design Rationale}

The computational domain is a straight circular main pipe with a single
circular branch intersecting it.  Two key length scales are fixed by the
literature:

\begin{itemize}
  \item \textbf{Main-pipe length} $= 20\,d_1 = 9200$\,mm.
        A minimum of 10\,$d_1$ upstream of the junction is provided
        as a practical compromise for turbulent profile development at
        $\mathrm{Re} \sim 10^7$~\cite{Schlichting2017}, and 10\,$d_1$ downstream
        provides sufficient mixing length for concentration measurements.
        Eames et al.\ showed that stratification persists for at least
        8\,$d_1$ downstream~\cite{Eames2022}.
  \item \textbf{Branch length} must provide at least 10 branch diameters of
        straight run for flow development, plus enough extra length to clear
        the main pipe at the shallowest planned angle.
\end{itemize}


\section{Setup}

Create one Workbench system per injection angle:

\begin{enumerate}
  \item In Workbench, drag a \textbf{Fluid Flow (Fluent)} system onto the
        project schematic.  Rename it \texttt{Theta\_090} (or whichever
        angle you are building first).
  \item Repeat for each of the remaining four angles, naming them
        \texttt{Theta\_030}, \texttt{Theta\_060}, \texttt{Theta\_120},
        and \texttt{Theta\_150}.
  \item Double-click the \textbf{Geometry} cell of the first system to
        open DesignModeler.
  \item Set \textbf{Units} to \textbf{Millimeter}.
\end{enumerate}

\noindent The steps below describe the geometry construction for a
\emph{single} base system.  Repeat the entire procedure for each of the
five angles, substituting the appropriate value of $\theta$.


\section{Main Pipeline}

\begin{enumerate}
  \item In the Tree Outline, click \textbf{XYPlane} $\rightarrow$
        \textbf{New Sketch} (creates \texttt{Sketch1}).
  \item Switch to the \textbf{Sketching} tab $\rightarrow$ \textbf{Draw}
        $\rightarrow$ \textbf{Circle}.  Click the origin and drag outward.
  \item \textbf{Dimensions} $\rightarrow$ \textbf{General}: click the circle
        edge and set the diameter to \textbf{460}.
  \item Switch back to \textbf{Modeling}, click \textbf{Extrude}:
  \begin{itemize}
    \item \textbf{Operation:} Add Material
    \item \textbf{Direction:} Normal
    \item \textbf{Extent Type:} Fixed
    \item \textbf{Depth:} \textbf{9200}
  \end{itemize}
  \item Click \textbf{Generate}.
\end{enumerate}

\noindent\fbox{\parbox{\dimexpr\textwidth-2\fboxsep-2\fboxrule}{%
\textbf{Result:} A 9200\,mm cylinder along the Z-axis representing the
main pipeline fluid domain.}}


\section{Injection Plane (Fixed per Base)}
\label{sec:injplane}

This plane sets the injection angle $\theta$ for this base system.
Because each angle has its own Workbench system, $\theta$ is entered as
a \textbf{fixed value}---it is \emph{not} parameterised.

\begin{enumerate}
  \item In the Tree Outline, click \textbf{ZXPlane} $\rightarrow$
        \textbf{New Plane}.
  \item In the Details View, apply:
  \begin{itemize}
    \item \textbf{Transform 1 (RMB):} \textbf{Offset Z} $= 4600$.
    \item \textbf{Transform 2 (RMB):} \textbf{Rotate about X} $= \theta$.
          Enter the fixed angle for this base system (e.g.\ $90^\circ$
          for \texttt{Theta\_090}, $30^\circ$ for \texttt{Theta\_030},
          etc.).
  \end{itemize}
  \item \textbf{Do not parameterise} this value.  Leave the checkbox
        blank.
  \item Click \textbf{Generate}.
\end{enumerate}

\medskip
\noindent\fbox{\parbox{\dimexpr\textwidth-2\fboxsep-2\fboxrule}{%
\textbf{Why 4600\,mm?}  The branch is placed at 10\,$d_1$ from the main
inlet, providing sufficient development length for the turbulent profile
to establish before reaching the junction.

\textbf{Why not parameterise $\theta$?}  Changing the injection angle
alters the branch--pipe intersection topology (face count, edge splits),
which breaks Named Selections and may cause mesh failures.  By fixing
$\theta$ per base and only sweeping $\beta$ and $R_v$, the geometry
topology is guaranteed to remain constant across all design points
within each base.}}


\section{Parameterised Branch Pipe}

\begin{enumerate}
  \item Select the new plane in the tree $\rightarrow$ \textbf{New Sketch}
        (creates \texttt{Sketch2}).
  \item \textbf{Sketching} $\rightarrow$ \textbf{Draw} $\rightarrow$
        \textbf{Circle}.  Click the origin of the new plane and draw a circle.
  \item \textbf{Dimensions} $\rightarrow$ \textbf{General}: set diameter to
        \textbf{115} (the $\beta = 0.25$ baseline, i.e.\ $0.25 \times 460$).
  \item \textbf{Parameterise} this dimension: click the blank square beside
        the value.  Name the parameter \texttt{Branch\_Dia}.
  \item Switch to \textbf{Modeling}, select \texttt{Sketch2}, click
        \textbf{Extrude}:
  \begin{itemize}
    \item \textbf{Operation:} \textbf{Add Material}
    \item \textbf{Direction:} Normal
    \item \textbf{Depth:} Enter the value from Table~\ref{tab:branchdepth}
          for this base system's angle.
  \end{itemize}
  \item \textbf{Parameterise} the depth: click the blank square beside the
        value.  Name the parameter \texttt{Branch\_Extrude\_Depth}.
  \item Click \textbf{Generate}.
\end{enumerate}

\subsection{Branch Extrude Depth Expression}

Because $\theta$ is fixed per base, the clearance term
$d_1/(2\sin\theta)$ is a known constant for each system.  In the
Workbench \textbf{Parameter Set}, define
\texttt{Branch\_Extrude\_Depth} with the expression:

\begin{equation}
  \texttt{Branch\_Extrude\_Depth}
  = 10 \times \texttt{Branch\_Dia}
  + \underbrace{\frac{d_1 / 2}{\sin\theta}}_{\text{constant per base}}
  \label{eq:branchdepth}
\end{equation}

The clearance constants for each base are pre-computed in
Table~\ref{tab:branchdepth}.

\begin{table}[H]
\centering
\caption{Branch extrude depth: fixed clearance term per base system.}
\label{tab:branchdepth}
\begin{tabular}{@{}lcll@{}}
\toprule
\textbf{Base} & $\theta$ & $d_1/(2\sin\theta)$ &
\textbf{Workbench Expression} \\
\midrule
\texttt{Theta\_030} & $30^\circ$  & 460\,mm &
\texttt{10 * Branch\_Dia + 460} \\
\texttt{Theta\_060} & $60^\circ$  & 266\,mm &
\texttt{10 * Branch\_Dia + 266} \\
\texttt{Theta\_090} & $90^\circ$  & 230\,mm &
\texttt{10 * Branch\_Dia + 230} \\
\texttt{Theta\_120} & $120^\circ$ & 266\,mm &
\texttt{10 * Branch\_Dia + 266} \\
\texttt{Theta\_150} & $150^\circ$ & 460\,mm &
\texttt{10 * Branch\_Dia + 460} \\
\bottomrule
\end{tabular}
\end{table}

\noindent This expression contains no trigonometric functions, eliminating
the degrees-vs-radians ambiguity that would arise if $\theta$ were
parameterised.


\section{Named Selections (in DesignModeler)}
\label{sec:namedsel}

Named Selections \textbf{must} be created inside DesignModeler rather than
in the Meshing application.  DesignModeler-based Named Selections are tied
to the \emph{generating feature} (e.g.\ the end face of Extrude\,1), so
they survive parametric geometry updates when \texttt{Branch\_Dia} changes.

Because $\theta$ is fixed per base system, the branch--pipe intersection
topology never changes within a single system---only the branch diameter
varies.  This guarantees that the face structure (and therefore the Named
Selections) remains stable across all 16 design points in each base.

\subsection{Enable Named Selection Transfer in Workbench}

Before proceeding, enable the transfer of Named Selections from
DesignModeler to Meshing and Fluent:

\begin{enumerate}
  \item In Workbench, go to \textbf{Tools} $\rightarrow$
        \textbf{Options} $\rightarrow$ \textbf{Geometry Import}.
  \item Ensure the \textbf{Named Selections} checkbox is enabled.
  \item Click \textbf{OK}.
\end{enumerate}

\subsection{Creating the Named Selections}

\subsubsection{Body-Level Named Selections}

These are used to scope mesh controls (Method, Sizing, Inflation) and
must be created \emph{before} the face-level boundary selections.

\begin{enumerate}
  \item In DesignModeler, go to \textbf{Tools} $\rightarrow$
        \textbf{Named Selection}.
  \item \textbf{\texttt{Fluid\_Domain}:} Select the
        \texttt{T\_Junction\_Fluid} body (the main pipe + branch solid).
        Right-click $\rightarrow$ \textbf{Create Named Selection}.
        Name it \texttt{Fluid\_Domain}.
  \item \textbf{\texttt{BOI\_Body}:} Select the \texttt{BOI\_Zone} body
        (the frozen sphere).  Create Named Selection $\rightarrow$
        \texttt{BOI\_Body}.
  \item Click \textbf{Generate}.
\end{enumerate}

\subsubsection{Face-Level Named Selections (Boundary Conditions)}

\begin{enumerate}
  \item With the complete geometry visible (main pipe + branch + BOI
        sphere), go to \textbf{Tools} $\rightarrow$
        \textbf{Named Selection}.
  \item \textbf{\texttt{main\_inlet}:} Select the circular face at
        $Z = 0$ (the flat end of the first extrusion).  Right-click
        $\rightarrow$ \textbf{Create Named Selection}.  Name it
        \texttt{main\_inlet}.
  \item \textbf{\texttt{outlet}:} Select the circular face at
        $Z = 9200$\,mm (the opposite end of the first extrusion).
        Create Named Selection $\rightarrow$ \texttt{outlet}.
  \item \textbf{\texttt{branch\_inlet}:} Select the flat end face of
        the branch extrusion (the face farthest from the main pipe).
        Create Named Selection $\rightarrow$ \texttt{branch\_inlet}.
  \item \textbf{\texttt{wall}:} Click \textbf{Select All} (to highlight
        every face), then hold \textbf{Ctrl} and click to deselect the
        three faces already named (\texttt{main\_inlet},
        \texttt{outlet}, \texttt{branch\_inlet}).  Create Named
        Selection $\rightarrow$ \texttt{wall}.
  \item Click \textbf{Generate}.
\end{enumerate}

\subsection{Post-Update Verification}

After the first design-point update within a base, open the
\textbf{Mesh} cell and check the \textbf{Named Selections} branch
in the Outline tree:

\begin{itemize}
  \item If all four selections display normally, the transfer is
        intact and the remaining design points will also succeed.
  \item If any selection shows a \textbf{yellow warning icon},
        return to DesignModeler, repair the affected Named Selection
        by re-picking the correct face, and regenerate.
\end{itemize}

\medskip
\noindent\fbox{\parbox{\dimexpr\textwidth-2\fboxsep-2\fboxrule}{%
\textbf{Why this matters:} If Named Selections are invalid, boundary
conditions in Fluent will map to the wrong faces or disappear entirely.
The five-base architecture eliminates the highest-risk topology change
(varying $\theta$), but verifying after the first $\beta$ update within
each base is still recommended.}}


\section{Material Assignment and Cleanup}

\begin{enumerate}
  \item Expand \textbf{1 Part, 1 Body} in the Tree Outline.  (The
        \texttt{BOI\_Zone} will appear as a second body under a
        \textbf{2 Parts, 2 Bodies} folder.)
  \item Right-click the pipe Solid $\rightarrow$ \textbf{Rename} to
        \texttt{T\_Junction\_Fluid}.
  \item Click \texttt{T\_Junction\_Fluid}.  In the Details View, change
        \textbf{Fluid/Solid} from \texttt{Solid} to \textbf{Fluid}.
  \item Verify that \texttt{BOI\_Zone} is also set to \textbf{Fluid}.
  \item Save the project and close DesignModeler.
\end{enumerate}


% ════════════════════════════════════════════════════════════════
%  CHAPTER 3 — Parameterised Meshing
% ════════════════════════════════════════════════════════════════
\chapter{Parameterised Mesh Generation}
\label{ch:mesh}

\section{Pre-Requisite: Body of Influence in DesignModeler}

A Body of Influence (BOI) provides localised mesh refinement in the mixing
zone without refining the entire domain.  Before entering the Meshing
application, create this zone in DesignModeler:

\begin{enumerate}
  \item \textbf{Create} $\rightarrow$ \textbf{Primitives} $\rightarrow$
        \textbf{Sphere}.
  \item Set \textbf{Origin Z} to \textbf{4600\,mm} (centred on the junction).
  \item Set \textbf{Operation} to \textbf{Add Frozen} (prevents merging with
        the pipe body).
  \item Set the baseline \textbf{Radius} to \textbf{690\,mm} ($= 1.5\,d_1$).
  \item \textbf{Parameterise} the radius: click the blank square beside the
        value and name the parameter \texttt{BOI\_Radius}.
  \item Click \textbf{Generate}.  Rename the new body to \texttt{BOI\_Zone}
        and set its material to \textbf{Fluid}.
\end{enumerate}

\subsection{BOI Radius Expression}

In the Parameter Set, define:

\begin{equation}
  \texttt{BOI\_Radius}
  = \max\!\bigl(1.5\,d_1,\;\; 8 \times \texttt{Branch\_Dia}\bigr)
  \label{eq:boiradius}
\end{equation}

\noindent\textbf{Justification:} Eames et al.~\cite{Eames2022} demonstrated
that hydrogen stratification persists for up to 8 main-pipe diameters
downstream of the injection point.  The refinement sphere must enclose this
region.  The $\max$ function ensures adequate coverage for both small and large
branch diameters.

\medskip
\noindent\fbox{\parbox{\dimexpr\textwidth-2\fboxsep-2\fboxrule}{%
\textbf{Workbench expression syntax:}
If \texttt{Branch\_Dia} is parameter \texttt{P2} and $d_1 = 460$\,mm, type:\\
\texttt{max(690, P2 * 8)}}}


\section{Global Mesh Settings}

\begin{enumerate}
  \item In Workbench, double-click the \textbf{Mesh} cell.
  \item Click \textbf{Mesh} in the Outline tree.  In the Details panel:
  \begin{itemize}
    \item \textbf{Physics Preference:} CFD
    \item \textbf{Solver Preference:} Fluent
  \end{itemize}
  \item Under \textbf{Sizing}:
  \begin{itemize}
    \item \textbf{Element Size:} 15\,mm \quad ($\approx d_1/30$; provides a
          coarse far-field mesh that keeps cell count manageable)
    \item \textbf{Growth Rate:} 1.15 \quad (industry standard for smooth
          transitions~\cite{ANSYSFluentUG2024})
    \item \textbf{Capture Curvature:} Yes
    \item \textbf{Curvature Normal Angle:} $15^\circ$ \quad (ensures circular
          cross-sections are geometrically faithful)
  \end{itemize}
  \item Under \textbf{Advanced}: set \textbf{Element Order} to \textbf{Linear}.
\end{enumerate}


\section{Named Selections}

Six Named Selections were created in DesignModeler
(Section~\ref{sec:namedsel}): two body-level (\texttt{Fluid\_Domain},
\texttt{BOI\_Body}) and four face-level (\texttt{main\_inlet},
\texttt{branch\_inlet}, \texttt{outlet}, \texttt{wall}).  Verify that
all six appear in the \textbf{Named Selections} branch of the Meshing
Outline tree.  Do \textbf{not} redefine them here---doing so would
create fragile, ID-based selections that break during parametric updates.

\medskip
\noindent\fbox{\parbox{\dimexpr\textwidth-2\fboxsep-2\fboxrule}{%
\textbf{Why body-level Named Selections matter:} All mesh controls
below (Method, BOI Sizing, Inflation) are scoped to
\texttt{Fluid\_Domain} or \texttt{BOI\_Body} via Named Selections
instead of direct geometry picking.  Direct picking (e.g.\ Ctrl+B on a
body) creates a hard-coded reference to an internal body ID.  When
\texttt{Branch\_Dia} changes and the geometry regenerates, these IDs
can shift, causing the mesh control to lose its scoping attachment and
silently fail.  Named-Selection-based scoping is resolved by
\emph{name}, making it immune to ID renumbering.}}


\section{Mesh Method}

\begin{enumerate}
  \item Right-click \textbf{Mesh} $\rightarrow$ \textbf{Insert}
        $\rightarrow$ \textbf{Method}.
  \item \textbf{Scoping Method:} \textbf{Named Selection}.
  \item \textbf{Named Selection:} \texttt{Fluid\_Domain}.
  \item \textbf{Method:} \textbf{MultiZone}.
  \item \textbf{Src/Trg Selection:} Automatic.
  \item \textbf{Free Mesh Type:} \textbf{Hexa Core}.
\end{enumerate}

\noindent\fbox{\parbox{\dimexpr\textwidth-2\fboxsep-2\fboxrule}{%
\textbf{Mesh method selection per base:} Because $\theta$ is fixed per
base, you can choose the optimal mesh method for each angle independently.
MultiZone with Hexa Core works well for $\theta \geq 60^\circ$.  For the
\texttt{Theta\_030} base (shallow intersection), MultiZone may fail due
to the complex junction topology.  In that case, switch the Method to
\textbf{Tetrahedrons} and convert to polyhedra inside Fluent
(\texttt{Mesh} $\rightarrow$ \texttt{Polyhedra} $\rightarrow$
\texttt{Convert}).  Polyhedral cells reduce cell count by $\sim$4$\times$
while maintaining accuracy~\cite{ANSYSFluentUG2024}.}}


\section{Body-of-Influence Sizing}

\begin{enumerate}
  \item Right-click \textbf{Mesh} $\rightarrow$ \textbf{Insert}
        $\rightarrow$ \textbf{Sizing}.
  \item \textbf{Scoping Method:} \textbf{Named Selection}.
  \item \textbf{Named Selection:} \texttt{Fluid\_Domain}.
  \item \textbf{Type:} \textbf{Body of Influence}.
  \item \textbf{Body of Influence:} Change \textbf{Scoping Method} to
        \textbf{Named Selection}, then select \texttt{BOI\_Body}.
  \item \textbf{Element Size:} \textbf{5\,mm}.
\end{enumerate}

\noindent\textbf{Justification:} For the baseline case ($\beta = 0.25$,
$d_2 = 115$\,mm), this gives $\sim$23 cells across the branch diameter,
sufficient to resolve the jet shear layer.  Gritskevich et
al.~\cite{Gritskevich2012} demonstrated adequate resolution of T-junction
mixing with comparable cell-per-diameter ratios in their OECD benchmark.


\section{Inflation Layer (Boundary Layer Mesh)}
\label{sec:inflation}

A resolved boundary layer ($y^+ \approx 1$) is mandatory for the
Spalart--Allmaras DDES model, which does not use wall
functions~\cite{SpalartAllmaras1994, Spalart2006}.

\subsection{First-Layer Height Derivation}

The first cell height $\Delta y_1$ for a target $y^+ = 1$ is derived as
follows~\cite{Schlichting2017}:

\begin{enumerate}
  \item \textbf{Reynolds number} based on main-pipe diameter and
        methane properties at $p = 6.9$\,MPa, $T = 284$\,K:
  \begin{align}
    \rho_1 &= \frac{M_{\mathrm{CH_4}}\, p}{R_g\, T_1}
            = \frac{0.01604 \times 6.9 \times 10^6}{8.314 \times 284}
            \approx 46.9\;\text{kg\,m}^{-3}
    \label{eq:rho1} \\[4pt]
    \mu_1  &\approx 1.25 \times 10^{-5}\;\text{Pa\,s}
            \qquad\text{(NIST~\cite{NIST2023} at 6.9\,MPa, 284\,K)}
    \label{eq:mu1} \\[4pt]
    \mathrm{Re}_{d_1}
           &= \frac{\rho_1\, u_1\, d_1}{\mu_1}
            = \frac{46.9 \times 10 \times 0.46}{1.25 \times 10^{-5}}
            \approx 1.73 \times 10^{7}
    \label{eq:Re}
  \end{align}

  \item \textbf{Skin-friction coefficient} via the Schlichting
        correlation for turbulent pipe flow~\cite{Schlichting2017}:
  \begin{equation}
    C_f = 0.058\;\mathrm{Re}_{d_1}^{\;-0.2}
        = 0.058 \times (1.73 \times 10^7)^{-0.2}
        \approx 2.11 \times 10^{-3}
    \label{eq:Cf}
  \end{equation}

  \item \textbf{Wall shear stress} and \textbf{friction velocity}:
  \begin{align}
    \tau_w &= \tfrac{1}{2}\, C_f\, \rho_1\, u_1^2
            = \tfrac{1}{2} \times 2.11 \times 10^{-3} \times 46.9 \times 100
            \approx 4.95\;\text{Pa}
    \label{eq:tauw} \\[4pt]
    u_\tau &= \sqrt{\tau_w / \rho_1}
            = \sqrt{4.95 / 46.9}
            \approx 0.325\;\text{m\,s}^{-1}
    \label{eq:utau}
  \end{align}

  \item \textbf{First cell height} for $y^+ = 1$:
  \begin{equation}
    \Delta y_1 = \frac{y^+\, \mu_1}{\rho_1\, u_\tau}
               = \frac{1 \times 1.25 \times 10^{-5}}{46.9 \times 0.325}
               \approx 8.2 \times 10^{-7}\;\text{m}
               \approx \boxed{0.001\;\text{mm}}
    \label{eq:dy1}
  \end{equation}
\end{enumerate}

\medskip
\noindent\fbox{\parbox{\dimexpr\textwidth-2\fboxsep-2\fboxrule}{%
\textbf{Result:} At 6.9\,MPa operating pressure, a first-layer height of
\textbf{0.001\,mm (1\,$\mu$m)} achieves $y^+ \approx 1$ for the methane
main flow.  The high gas density at this pressure ($\sim$47\,kg\,m$^{-3}$
versus $\sim$0.66\,kg\,m$^{-3}$ at atmospheric) is the dominant factor
requiring such a thin first cell.  This value should be verified against
the post-processed $y^+$ field after the first simulation run.}}

\subsection{Inflation Settings}

\begin{enumerate}
  \item Right-click \textbf{Mesh} $\rightarrow$ \textbf{Insert}
        $\rightarrow$ \textbf{Inflation}.
  \item \textbf{Scoping Method:} \textbf{Named Selection}.
  \item \textbf{Named Selection:} \texttt{Fluid\_Domain}.
  \item \textbf{Boundary Scoping Method:} \textbf{Named Selections}.
  \item \textbf{Boundary:} Select the \texttt{wall} named selection.
  \item \textbf{Inflation Option:} \textbf{First Layer Thickness}.
  \item \textbf{First Layer Height:} \textbf{0.001\,mm}.
  \item \textbf{Maximum Layers:} \textbf{20}.
  \item \textbf{Growth Rate:} \textbf{1.2}.
\end{enumerate}

\noindent\textbf{Layer count justification:} With 20 layers and a growth
rate of 1.2, the inflation stack spans approximately
$\Delta y_1 \times (1.2^{20} - 1)/(1.2 - 1) \approx 0.001 \times 186.7
\approx 0.19$\,mm from the wall.  The turbulent boundary-layer thickness
$\delta$ for the main pipe can be estimated as
$\delta \approx 0.37\,d_1\,\mathrm{Re}_{d_1}^{-0.2}
\approx 0.37 \times 460 \times (1.73 \times 10^7)^{-0.2}
\approx 6.2$\,mm~\cite{Schlichting2017}.
The inflation stack covers the viscous sublayer and buffer layer; the
outer boundary layer is resolved by the BOI cells (5\,mm).


\section{Mesh Quality Targets}

After generating the mesh (\textbf{Mesh} $\rightarrow$ \textbf{Generate}),
verify the following quality metrics in the \textbf{Statistics} panel:

\begin{table}[H]
\centering
\caption{Mesh quality acceptance criteria.}
\label{tab:meshqual}
\begin{tabular}{@{}lll@{}}
\toprule
\textbf{Metric} & \textbf{Target} & \textbf{Reference} \\
\midrule
Minimum Orthogonal Quality & $> 0.10$ & \cite{ANSYSFluentUG2024} \\
Maximum Skewness           & $< 0.95$ & \cite{ANSYSFluentUG2024} \\
Maximum Aspect Ratio       & $< 100$  & \cite{ANSYSFluentUG2024} \\
\bottomrule
\end{tabular}
\end{table}


% ════════════════════════════════════════════════════════════════
%  CHAPTER 4 — Fluent Setup: General and Models
% ════════════════════════════════════════════════════════════════
\chapter{Fluent Setup --- General and Physical Models}
\label{ch:models}

Double-click the \textbf{Setup} cell in Workbench to launch Fluent.
When prompted, select \textbf{Double Precision} and set the number of
\textbf{Solver Processes} to 8 (appropriate for a $\sim$3\,M cell mesh on an
i7-14700 with 8 P-cores).

\section{General Settings}

Navigate to \textbf{Setup} $\rightarrow$ \textbf{General}:

\begin{enumerate}
  \item \textbf{Type:} Pressure-Based.  \quad(Mach $\approx 0.02$)
  \item \textbf{Velocity Formulation:} Absolute.
  \item \textbf{Time:} \textbf{Transient}.  \quad(SA-DDES is inherently
        unsteady; no steady RANS precursor.)
  \item \textbf{Gravity:} \textbf{Enabled}.  Set $g_x = 0$,
        $g_y = -9.81$, $g_z = 0$.  \quad(Buoyancy drives hydrogen
        stratification; density ratio CH$_4$:H$_2$ $\approx 8$ at
        6.9\,MPa.)
\end{enumerate}

\section{Viscous Model (Turbulence)}

Navigate to \textbf{Models} $\rightarrow$ \textbf{Viscous}:

\begin{enumerate}
  \item \textbf{Model:} \textbf{Spalart--Allmaras (1\,eqn)}.
  \item Under \textbf{Spalart--Allmaras Options}:
  \begin{itemize}
    \item \textbf{Production:} Change from Vorticity-Based (default) to
          \textbf{Strain/Vorticity-Based}.
    \item \textbf{Curvature Correction:} \textbf{Enable}.
  \end{itemize}
  \item In the left list, select \textbf{Detached Eddy Simulation (DES)}.
  \item \textbf{DES Option:} \textbf{Delayed DES (DDES)} (default).
  \item \textbf{Model Constants:} Leave at defaults (well-validated).
  \item Click \textbf{OK}.
\end{enumerate}

\noindent\textbf{Rationale:}
\begin{itemize}
  \item SA-DDES chosen by Eames et al.~\cite{Eames2022}; DDES shielding
        prevents premature LES activation in the boundary
        layer~\cite{Spalart2006, Shur2008}.
  \item \textbf{Strain/Vorticity-Based Production} reduces over-prediction
        of eddy viscosity in high-shear regions such as the hydrogen jet
        mixing zone.  This is a known improvement for free shear
        flows~\cite{SpalartAllmaras1994}.
  \item \textbf{Curvature Correction} prevents the SA model from
        artificially damping resolved vortices in the pipe curvature and
        swirling wake downstream of the junction~\cite{Spalart2006}.
\end{itemize}

\medskip
\noindent\fbox{\parbox{\dimexpr\textwidth-2\fboxsep-2\fboxrule}{%
\textbf{Visual confirmation:} After applying settings, the Viscous Model
dialog should display:\\[4pt]
\begin{tabular}{@{}ll@{}}
Model:               & Detached Eddy Simulation (DES) \\
RANS Model:          & Spalart--Allmaras \\
DES Option:          & Delayed DES (DDES) \\
Production:          & Strain/Vorticity-Based \\
Curvature Correction:& Enabled \\
\end{tabular}}}

\section{Species Transport}

Navigate to \textbf{Models} $\rightarrow$ \textbf{Species}:

\begin{enumerate}
  \item Select \textbf{Species Transport}.
  \item \textbf{Reactions:} Disabled.  \quad(Mechanical mixing study, not
        combustion; saves significant computation.)
  \item \textbf{Diffusion Energy Source:} Enabled.  \quad(H$_2$ and CH$_4$
        have very different $c_p$ and thermal conductivities; as they
        diffuse into each other, enthalpy is transported.)
  \item \textbf{Full Multicomponent Diffusion:} Disabled.  \quad(With only
        two species---H$_2$ and CH$_4$---Fickian binary diffusion is
        exact~\cite{Poling2001}.  Enable only if modelling natural gas as
        a multi-component blend.)
\end{enumerate}

\section{Energy Equation}

\textbf{Enable} the Energy Equation.  At 6.9\,MPa the gas is highly
compressible; the solver needs the energy equation to recalculate
density via the ideal gas law ($\rho = p\,M / R_g\,T$) as pressure and
mixing ratios fluctuate in the T-junction~\cite{ANSYSFluent2024}.


% ════════════════════════════════════════════════════════════════
%  CHAPTER 5 — Materials and Operating Conditions
% ════════════════════════════════════════════════════════════════
\chapter{Materials and Operating Conditions}
\label{ch:materials}

\section{Phase 1 --- Import the Pure Gases}

Fluent has separate databases for \textbf{Mixtures} and pure
\textbf{Fluids}.  The pure species must be copied into your project
before a custom mixture can reference them.

\begin{enumerate}
  \item In the Fluent tree, expand \textbf{Materials}.
  \item Double-click \textbf{Fluid} (default entry is \texttt{air}).
        The \textit{Create/Edit Materials} window opens.
  \item Click \textbf{Fluent Database\ldots} on the right.
  \item In the database window, ensure \textbf{Material Type} is set to
        \textbf{\texttt{fluid}} (not \texttt{mixture}---selecting
        \texttt{mixture} shows pre-built templates like
        \texttt{methane-air}, which is not what we want).
  \item Scroll to \textbf{methane (ch4)}, click it, then click
        \textbf{Copy}.
  \item Scroll to \textbf{hydrogen (h2)}, click it, then click
        \textbf{Copy}.
  \item Click \textbf{Close}, then close the \textit{Create/Edit
        Materials} window.
\end{enumerate}

\section{Phase 2 --- Build the Custom Mixture}

\begin{enumerate}
  \item Under \textbf{Materials}, double-click \textbf{Mixture}
        (\texttt{mixture-template}).
  \item In the row \textbf{Mixture Species}, click \textbf{Edit\ldots}.
  \item The \textit{Species} dialog opens with \textit{Available
        Materials} (left) and \textit{Selected Species} (right).
  \item Add \textbf{ch4} and \textbf{h2} from left to right.
  \item Remove \textbf{o2} and \textbf{n2} from the right list.
  \item \textbf{Order matters:} the last species in the right list is
        the ``balance'' species (computed from $1 - \sum Y_i$).  Since
        the pipe is mostly CH$_4$, ensure the list reads:
        \begin{enumerate}
          \item \texttt{h2}
          \item \texttt{ch4} \quad(bottom = balance species)
        \end{enumerate}
  \item Click \textbf{OK}.
\end{enumerate}

\section{Phase 3 --- Set Mixture Physical Properties}

Back on the \textit{Create/Edit Materials} window for the mixture,
set each property drop-down:

\begin{table}[H]
\centering
\caption{Mixture-level property settings.}
\label{tab:mixprops}
\begin{tabular}{@{}ll@{}}
\toprule
\textbf{Property} & \textbf{Setting} \\
\midrule
Density              & \textbf{ideal-gas} \\
$c_p$ (Specific Heat)& \textbf{mixing-law} \\
Thermal Conductivity & \textbf{ideal-gas-mixing-law} \\
Viscosity            & \textbf{ideal-gas-mixing-law} \\
Mass Diffusivity     & \textbf{kinetic-theory} \\
\bottomrule
\end{tabular}
\end{table}

\noindent Click \textbf{Change/Create}, then \textbf{Close}.

\medskip
\noindent\textbf{Note:} Ideal-gas is a modelling simplification
(Eames et al.\ used it~\cite{Eames2022}).  Using mixing-law and
ideal-gas-mixing-law for $c_p$, conductivity, and viscosity lets
Fluent weight each property by local species concentration
automatically.  Temperature-dependent polynomials are negligible for
$\Delta T < 10$\,K~\cite{Poling2001}.

\section{Individual Species Properties (Verification)}

After the mixture is created, verify the per-species constants match
NIST data~\cite{NIST2023} at 6.9\,MPa:

\begin{table}[H]
\centering
\caption{Species thermophysical properties.}
\label{tab:matprops}
\begin{tabular}{@{}llll@{}}
\toprule
\textbf{Property} & \textbf{CH$_4$} & \textbf{H$_2$} & \textbf{Source} \\
\midrule
Molecular weight (kg\,kmol$^{-1}$) & 16.043 & 2.016 & NIST \\
$c_p$ (J\,kg$^{-1}$\,K$^{-1}$)    & 2580 & 14\,500 & NIST \\
Thermal conductivity (W\,m$^{-1}$\,K$^{-1}$) & 0.038 & 0.185 & NIST \\
Dynamic viscosity (Pa\,s)           & $1.25 \times 10^{-5}$ & $9.0 \times 10^{-6}$ & NIST \\
\bottomrule
\end{tabular}
\end{table}

\noindent These values are loaded from the Fluent database when you
copy the pure fluids.  If any differ significantly from the table above,
edit them manually in the species material panel.

\section{Turbulent Schmidt Number}
\label{sec:sct}

Default $\mathrm{Sc}_t = 0.7$; for H$_2$ jets in crossflow,
$\mathrm{Sc}_t = 0.5$ improves mixing prediction.  Change via TUI:

\begin{lstlisting}[language={},numbers=none]
/define/models/viscous/turbulent-schmidt-number 0.5
\end{lstlisting}

\section{Operating Conditions}

Navigate to \textbf{Cell Zone Conditions} $\rightarrow$
\textbf{Operating Conditions}:

\begin{enumerate}
  \item \textbf{Operating Pressure:} Enter \textbf{6900000}
        (6.9\,MPa; do not use commas).
  \item \textbf{Reference Pressure Location:} Leave $X$, $Y$, $Z$ at
        default (0, 0, 0).
  \item \textbf{Gravity:} Check the \textbf{Gravity} box.  Set
        $g_y = -9.81$\,m\,s$^{-2}$ (branch points along $Y$ after
        the 90$^\circ$ ZX-plane rotation); leave $g_x = g_z = 0$.
  \item \textbf{Specified Operating Density:} Check the box, enter
        \textbf{0}.
\end{enumerate}

\medskip
\noindent\fbox{\parbox{\dimexpr\textwidth-2\fboxsep-2\fboxrule}{%
\textbf{Critical:} Incorrect operating pressure changes density by
two orders of magnitude, invalidating all physics.}}

\medskip
\noindent\fbox{\parbox{\dimexpr\textwidth-2\fboxsep-2\fboxrule}{%
\textbf{Why Operating Density = 0?}  With gravity enabled, Fluent
computes buoyancy as $(\rho - \rho_0)\,\mathbf{g}$.  For a
compressible, multi-species mixture (ideal gas with variable
composition), the local density $\rho$ varies significantly with
pressure, temperature, and species concentration.  A non-zero
$\rho_0$ introduces an artificial hydrostatic baseline that can
destabilise the calculation.  Setting $\rho_0 = 0$ reduces the
buoyancy term to $\rho\,\mathbf{g}$, using the true local density
--- the correct formulation for compressible
flows~\cite{ANSYSFluent2024}.}}


% ════════════════════════════════════════════════════════════════
%  CHAPTER 6 — Boundary Conditions
% ════════════════════════════════════════════════════════════════
\chapter{Boundary Conditions (Parameterised)}
\label{ch:bc}

Expand \textbf{Boundary Conditions} in the Fluent tree.  All named
selections from meshing appear automatically.

\section{Main Inlet (\texttt{main\_inlet}) --- The Natural Gas Crossflow}

Double-click \texttt{main\_inlet}.  Ensure \textbf{Type} is
\texttt{velocity-inlet}.

\begin{enumerate}
  \item \textbf{Momentum Tab:}
  \begin{itemize}
    \item Velocity Magnitude: \textbf{10} (m/s).
    \item Turbulence Specification Method: \textbf{Intensity and
          Hydraulic Diameter}.
    \item Turbulent Intensity: \textbf{5} (\%).
    \item Hydraulic Diameter: \textbf{0.46} (m).
  \end{itemize}
  \item \textbf{Thermal Tab:}
  \begin{itemize}
    \item Temperature: \textbf{284} (K).
  \end{itemize}
  \item \textbf{Species Tab:}
  \begin{itemize}
    \item Species Mass Fractions (h2): \textbf{0}.
          \quad(Pure methane.)
  \end{itemize}
  \item Click \textbf{Apply} and close.
\end{enumerate}

\section{Branch Inlet (\texttt{branch\_inlet}) --- The Hydrogen Jet}

Double-click \texttt{branch\_inlet}.  Ensure \textbf{Type} is
\texttt{velocity-inlet}.

\begin{enumerate}
  \item \textbf{Momentum Tab:}
  \begin{itemize}
    \item Velocity Magnitude: \textbf{32} (m/s).
    \item \textbf{Parameterise:} Click the small drop-down arrow (or
          blank square) to the right of the velocity value box
          $\rightarrow$ select \textbf{New Input Parameter\ldots}
          $\rightarrow$ name it \texttt{branch\_velocity} $\rightarrow$
          click OK.  A blue ``P'' appears next to the box, linking it to
          the Workbench parameter table
          (Chapter~\ref{ch:param}).
    \item Turbulence Specification Method: \textbf{Intensity and
          Hydraulic Diameter}.
    \item Turbulent Intensity: \textbf{5} (\%).
    \item Hydraulic Diameter: \textbf{0.115} (m).
  \end{itemize}
  \item \textbf{Thermal Tab:}
  \begin{itemize}
    \item Temperature: \textbf{293} (K).
  \end{itemize}
  \item \textbf{Species Tab:}
  \begin{itemize}
    \item Species Mass Fractions (h2): \textbf{1}.
          \quad(100\,\% pure hydrogen.)
  \end{itemize}
  \item Click \textbf{Apply} and close.
\end{enumerate}

\noindent\textbf{Velocity derivation:}
$\mathcal{D} = \beta^2 (u_2/u_1)$ $\;\Rightarrow\;$
$0.2 = (0.25)^2 (u_2/10)$ $\;\Rightarrow\;$ $u_2 = 32$\,m\,s$^{-1}$.
In the Parameter Set, use
$\texttt{branch\_velocity} = R_v \times 10$ to sweep velocity ratio.

\section{Outlet (\texttt{outlet}) --- The Downstream Exhaust}

Double-click \texttt{outlet}.  Ensure \textbf{Type} is
\texttt{pressure-outlet}.

\begin{enumerate}
  \item \textbf{Momentum Tab:}
  \begin{itemize}
    \item Gauge Pressure: \textbf{0}.  \quad(0 gauge = 6.9\,MPa
          absolute, since operating pressure is already set.)
    \item Backflow Turbulence Specification: \textbf{Intensity and
          Hydraulic Diameter}.
    \item Backflow Turbulent Intensity: \textbf{5} (\%).
    \item Backflow Hydraulic Diameter: \textbf{0.46} (m).
  \end{itemize}
  \item \textbf{Thermal Tab:}
  \begin{itemize}
    \item Backflow Total Temperature: \textbf{284} (K).
  \end{itemize}
  \item \textbf{Species Tab:}
  \begin{itemize}
    \item Backflow Species Mass Fractions (h2): \textbf{0}.
  \end{itemize}
  \item Click \textbf{Apply} and close.
\end{enumerate}

\noindent\textit{Backflow parameters are safety nets Fluent uses if the
swirling vortex momentarily pulls fluid backwards through the exit.}

\section{Walls (\texttt{wall}) --- The Pipe Boundaries}

Double-click \texttt{wall}.  Ensure \textbf{Type} is \texttt{wall}.

\begin{enumerate}
  \item \textbf{Momentum Tab:} Verify \textbf{Stationary Wall} and
        \textbf{No Slip} (defaults).
  \item \textbf{Thermal Tab:} Set thermal condition to \textbf{Heat
        Flux} with value \textbf{0}\,W\,m$^{-2}$ (adiabatic --- buried
        pipe, no heat exchange~\cite{Eames2022}).
  \item Click \textbf{Apply} and close.
\end{enumerate}


% ════════════════════════════════════════════════════════════════
%  CHAPTER 7 — Solution Methods and Controls
% ════════════════════════════════════════════════════════════════
\chapter{Solution Methods and Controls}
\label{ch:solution}

\section{Pressure--Velocity Coupling}

Navigate to \textbf{Solution} $\rightarrow$ \textbf{Methods}.

\begin{enumerate}
  \item Under \textbf{Pressure-Velocity Coupling}, click the
        \textbf{Scheme} drop-down.
  \item Select \textbf{PISO}.
  \item Two options appear below (Skewness Correction, Neighbor
        Correction) --- leave at defaults (usually~1).
\end{enumerate}

\section{Spatial Discretisation}

Still in \textbf{Solution} $\rightarrow$ \textbf{Methods}, set each
drop-down in the \textbf{Spatial Discretization} section:

\begin{enumerate}
  \item \textbf{Gradient:} \texttt{Least Squares Cell Based} (default).
  \item \textbf{Pressure:} Change to \textbf{\texttt{PRESTO!}}
        \quad(crucial for the buoyancy/gravity settings).
  \item \textbf{Momentum:} Change to \textbf{\texttt{Bounded Central
        Differencing}} \quad(the most important setting for DDES).
  \item \textbf{Modified Turbulent Viscosity:} Change to
        \textbf{\texttt{Second Order Upwind}}.
  \item \textbf{Species:} Change to
        \textbf{\texttt{Second Order Upwind}}.
  \item \textbf{Energy:} Change to
        \textbf{\texttt{Second Order Upwind}}.
\end{enumerate}

\medskip
\noindent\fbox{\parbox{\dimexpr\textwidth-2\fboxsep-2\fboxrule}{%
\textbf{Critical:} Standard Upwind schemes for momentum create
numerical friction that instantly kills the turbulent swirling eddies
DDES is designed to resolve.  Bounded Central Differencing is the
default DES setting in Fluent and must not be
changed~\cite{ANSYSFluent2024, Menter2012}.}}

\section{Transient Formulation}

Still in \textbf{Solution} $\rightarrow$ \textbf{Methods}, at the
bottom of the panel:

\begin{enumerate}
  \item Under \textbf{Transient Formulation}, click the drop-down.
  \item Select \textbf{\texttt{Bounded Second Order Implicit}}.
\end{enumerate}

\noindent Prevents oscillations while maintaining temporal
accuracy~\cite{Menter2012}.

\section{Under-Relaxation Factors}

Navigate to \textbf{Solution} $\rightarrow$ \textbf{Controls}.  These
act as mathematical ``brakes'' preventing divergence if pressure changes
too violently between time steps.  Match the following values:

\begin{table}[H]
\centering
\caption{Under-relaxation factors (transient PISO).}
\label{tab:urf}
\begin{tabular}{@{}ll@{}}
\toprule
\textbf{Variable} & \textbf{Factor} \\
\midrule
Pressure                    & 0.3 \\
Momentum                    & 0.7 \\
Modified Turbulent Viscosity & 0.8 \\
Species                     & 0.9 \\
Energy                      & 1.0 \\
\bottomrule
\end{tabular}
\end{table}

\noindent If \textbf{Density} or \textbf{Body Forces} appear in the
list, leave them at their default value (1.0).


% ════════════════════════════════════════════════════════════════
%  CHAPTER 8 — Run Setup, Monitors, and Execution
% ════════════════════════════════════════════════════════════════
\chapter{Run Setup, Monitors, and Execution}
\label{ch:run}

Follow these steps \textbf{in order} after setting the under-relaxation
factors.  Do not skip any.

% ── Step 1 ──────────────────────────────────────────────────────
\section{Step 1: Verify Solution Methods (Quick Check)}

Go to \textbf{Solution} $\rightarrow$ \textbf{Methods} and confirm:

\begin{itemize}
  \item Pressure--Velocity Coupling: \texttt{PISO}
  \item Transient Formulation: \texttt{Bounded Second Order Implicit}
  \item Momentum: \texttt{Bounded Central Differencing} (\emph{not}
        Upwind)
  \item Pressure: \texttt{PRESTO!}
  \item All others: \texttt{Second Order Upwind}
\end{itemize}

\noindent\fbox{\parbox{\dimexpr\textwidth-2\fboxsep-2\fboxrule}{%
If Momentum is \texttt{Second Order Upwind}, change it now ---
otherwise the DES will not resolve turbulence.}}

% ── Step 2 ──────────────────────────────────────────────────────
\section{Step 2: Run Calculation Parameters}

Go to \textbf{Solution} $\rightarrow$ \textbf{Run Calculation} and
enter:

\begin{table}[H]
\centering
\caption{Run calculation settings.}
\label{tab:transient}
\begin{tabular}{@{}lll@{}}
\toprule
\textbf{Field} & \textbf{Value} & \textbf{Explanation} \\
\midrule
Time Step Size (s)          & \texttt{2e-4}  & $\Delta t = 2\times10^{-4}$\,s; CFL $\approx 1.3$ \\
Number of Time Steps        & \texttt{25000} & $0.0002 \times 25{,}000 = 5$\,s total \\
Max Iterations / Time Step  & \texttt{20}    & Enough for residuals to drop 1--2 orders \\
\bottomrule
\end{tabular}
\end{table}

\noindent\textbf{Do not click Calculate yet} --- first set up monitors
and autosave.

% ── Step 3 ──────────────────────────────────────────────────────
\section{Step 3: Residual Monitors}

\begin{enumerate}
  \item \textbf{Solution} $\rightarrow$ \textbf{Monitors}
        $\rightarrow$ \textbf{Residuals} $\rightarrow$ \textbf{Edit}.
  \item Check all equations: continuity, x-velocity, y-velocity,
        z-velocity, modified turbulent viscosity (\texttt{nut}),
        species (\texttt{h2}), energy.
  \item \textbf{Convergence Criterion:} \texttt{1e-4} for all; set
        energy to \texttt{1e-6}.
  \item Ensure \textbf{Plot} is enabled $\rightarrow$ click \textbf{OK}.
\end{enumerate}

% ── Step 4 ──────────────────────────────────────────────────────
\section{Step 4: Surface Monitors (Track Mixing)}

Right-click \textbf{Monitors} $\rightarrow$ \textbf{New} $\rightarrow$
\textbf{Surface Monitor\ldots} for each of the following:

\medskip
\noindent\textbf{Monitor 1 --- Outlet H$_2$ mass fraction:}
\begin{itemize}
  \item Name: \texttt{outlet\_h2}
  \item Report Type: \texttt{Mass-Weighted Average}
  \item Field Variable: Species $\rightarrow$ Mass Fraction of h2
  \item Surfaces: \texttt{outlet}
  \item Click \textbf{Create}.
\end{itemize}

\noindent\textbf{Monitor 2 --- Outlet mass flow (mass balance):}
\begin{itemize}
  \item Name: \texttt{outlet\_massflow}
  \item Report Type: \texttt{Mass Flow Rate}
  \item Surfaces: \texttt{outlet}
  \item Click \textbf{Create}.
\end{itemize}

\noindent\textbf{Monitor 3 --- Main inlet mass flow:}
\begin{itemize}
  \item Name: \texttt{main\_inlet\_massflow}
  \item Report Type: \texttt{Mass Flow Rate}
  \item Surfaces: \texttt{main\_inlet}
  \item Click \textbf{Create}.
\end{itemize}

\noindent(Repeat for \texttt{branch\_inlet} if desired.)

% ── Step 5 ──────────────────────────────────────────────────────
\section{Step 5: Point Monitors (Wall Concentrations)}

Track H$_2$ mass fraction at the pipe-top wall ($y = +0.23$\,m) at
four downstream locations.  Right-click \textbf{Monitors} $\rightarrow$
\textbf{New} $\rightarrow$ \textbf{Point Monitor\ldots} for each row:

\begin{table}[H]
\centering
\caption{Point monitor coordinates (junction at $Z = 4.6$\,m,
         $d_1 = 0.46$\,m).}
\label{tab:pointmon}
\begin{tabular}{@{}lcccc@{}}
\toprule
\textbf{Name} & $X$ (m) & $Y$ (m) & $Z$ (m) & \textbf{Position} \\
\midrule
\texttt{wall\_h2\_z2d} & 0 & 0.23 & 5.52 & $z/d_1 = 2$ \\
\texttt{wall\_h2\_z4d} & 0 & 0.23 & 6.44 & $z/d_1 = 4$ \\
\texttt{wall\_h2\_z6d} & 0 & 0.23 & 7.36 & $z/d_1 = 6$ \\
\texttt{wall\_h2\_z8d} & 0 & 0.23 & 8.28 & $z/d_1 = 8$ \\
\bottomrule
\end{tabular}
\end{table}

\noindent For each: \textbf{Variable} $\rightarrow$ Species
$\rightarrow$ Mass Fraction of h2 $\rightarrow$ click \textbf{Create}.

% ── Step 6 ──────────────────────────────────────────────────────
\section{Step 6: Autosave (Prevents Data Loss)}

\begin{enumerate}
  \item Go to \textbf{Solution} $\rightarrow$ \textbf{Run Calculation}
        $\rightarrow$ click \textbf{Activities} (bottom of panel).
  \item Click \textbf{Create} $\rightarrow$ \textbf{Autosave}.
  \item Save Data File Every (Time Steps): \textbf{500}.
  \item Save Case File Every: \textbf{500}.
  \item \textit{Retain Only the Most Recent Files:} Optional (uncheck
        to keep all).
  \item Click \textbf{OK}, then close Activities.
\end{enumerate}

\noindent This saves every 0.1\,s ($500 \times 2\times10^{-4}$).
Adjust if disk space is limited.

% ── Step 7 ──────────────────────────────────────────────────────
\section{Step 7: Initialisation}

\begin{enumerate}
  \item Navigate to \textbf{Solution} $\rightarrow$
        \textbf{Initialization}.
  \item Select \textbf{Hybrid Initialization} $\rightarrow$ click
        \textbf{Initialize}.
  \item \textbf{Patch branch region:}
  \begin{itemize}
    \item Click \textbf{Patch\ldots}.
    \item Create a cell register (sphere/cylinder around the branch
          junction).
    \item Set H$_2$ mass fraction $= 1.0$ and velocity $= 32$\,m\,s$^{-1}$
          along the branch axis.
  \end{itemize}
\end{enumerate}

% ── Step 8 ──────────────────────────────────────────────────────
\section{Step 8: Start the Calculation}

Return to \textbf{Solution} $\rightarrow$ \textbf{Run Calculation}
$\rightarrow$ click \textbf{Calculate}.

Residuals and monitor plots will update in real time.

% ── Step 9 ──────────────────────────────────────────────────────
\section{Step 9: During the Run --- What to Watch}

\begin{table}[H]
\centering
\caption{Runtime diagnostics.}
\label{tab:diagnostics}
\begin{tabular}{@{}ll@{}}
\toprule
\textbf{Observation} & \textbf{Action} \\
\midrule
Residuals drop 1--2 orders each time step & Normal \\
Residuals stagnate or increase & Reduce $\Delta t$ to \texttt{1e-4}
  for a few hundred steps \\
Outlet H$_2$ mass fraction $\to$ $\sim$0.149 ($\mathcal{D}=0.2$) & Good \\
Mass flow imbalance $> 1$\,\% & Check boundary conditions \\
Temperature goes negative & Reduce $\Delta t$; check energy equation \\
\bottomrule
\end{tabular}
\end{table}

% ── Step 10 ─────────────────────────────────────────────────────
\section{Step 10: Enable Time Statistics (After $t = 4$\,s)}

After the first 20\,000 time steps ($t = 4$\,s), the initial
transients have washed out.  Enable time-averaged statistics for the
final 1\,s:

\begin{enumerate}
  \item \textbf{Solution} $\rightarrow$ \textbf{Run Calculation}
        $\rightarrow$ under \textbf{Data Sampling}, check
        \textbf{Enable}.
  \item \textbf{Sampling Interval:} \texttt{1} (every time step).
  \item Let the simulation continue to $t = 5$\,s (step 25\,000).
\end{enumerate}

\medskip
\noindent\fbox{\parbox{\dimexpr\textwidth-2\fboxsep-2\fboxrule}{%
\textbf{Summary checklist:}\\[4pt]
\begin{tabular}{@{}ll@{}}
Under-relaxation factors set              & $\square$ \\
Solution methods verified (BCD for momentum) & $\square$ \\
Run parameters entered ($\Delta t = 2\times10^{-4}$, 25\,000 steps, 20 iter) & $\square$ \\
Residual monitors enabled                  & $\square$ \\
Surface monitors created (outlet H$_2$, mass flows) & $\square$ \\
Point monitors created (wall H$_2$ at $2d$, $4d$, $6d$, $8d$) & $\square$ \\
Autosave configured (every 500 steps)      & $\square$ \\
Initialisation complete (Hybrid + Patch)   & $\square$ \\
Clicked \textbf{Calculate}                 & $\square$ \\
\end{tabular}}}


% ════════════════════════════════════════════════════════════════
%  CHAPTER 9 — Workbench Parameterisation
% ════════════════════════════════════════════════════════════════
\chapter{Workbench Parameterisation and Design Points}
\label{ch:param}

\section{Parameter Summary}

Each of the five base systems contains the same set of Workbench
parameters.  After completing the Geometry, Mesh, and Fluent setup for
a single base, the following parameters should be visible when you
double-click the \textbf{Parameter Set} bar:

\begin{table}[H]
\centering
\caption{Workbench parameters (identical in each of the five bases).}
\label{tab:allparams}
\begin{tabular}{@{}llll@{}}
\toprule
\textbf{Parameter Name} & \textbf{Source} & \textbf{Type} & \textbf{Baseline} \\
\midrule
\texttt{Branch\_Dia}            & DesignModeler (Sketch dim.)    & Input  & 115\,mm \\
\texttt{Branch\_Extrude\_Depth} & DesignModeler (Extrude depth)  & Output & (expression) \\
\texttt{BOI\_Radius}            & DesignModeler (Sphere radius)  & Output & (expression) \\
\texttt{branch\_velocity}       & Fluent (BC velocity mag.)      & Input  & 32\,m\,s$^{-1}$ \\
\bottomrule
\end{tabular}
\end{table}

\noindent \texttt{Branch\_Extrude\_Depth} and \texttt{BOI\_Radius} are
\emph{output} parameters driven by expressions linked to
\texttt{Branch\_Dia}.  They update automatically.  The injection angle
$\theta$ is \textbf{not} a Workbench parameter---it is baked into each
base system's geometry.


\section{Defining Expressions}

In the Parameter Set table, click on the \textbf{Expression} cell for
each driven parameter and enter:

\begin{enumerate}
  \item \texttt{Branch\_Extrude\_Depth}: Enter the expression from
        Table~\ref{tab:branchdepth} corresponding to this base's angle.

        \emph{Example for \texttt{Theta\_090}:}
        \texttt{10 * Branch\_Dia + 230}

  \item \texttt{BOI\_Radius}: \quad
        \texttt{max(690, Branch\_Dia * 8)}

        \emph{(690\,mm $= 1.5 \times d_1$; identical across all bases.)}
\end{enumerate}


\section{Creating the Design-Point Table}

Within each base system, the design-point table sweeps \texttt{Branch\_Dia}
($\beta$) and \texttt{branch\_velocity} ($R_v$).  With 4 levels of each,
the full factorial yields $4 \times 4 = 16$ design points per base.

\begin{enumerate}
  \item In the Parameter Set, click \textbf{Design Points} (bottom panel).
  \item Populate the table using the mapping:
  \begin{align}
    \texttt{Branch\_Dia} &= \beta \times 460 \notag \\
    \texttt{branch\_velocity} &= R_v \times 10 \notag
  \end{align}
  \item The complete 16-point table for each base is shown in
        Table~\ref{tab:designpoints}.
\end{enumerate}

\begin{table}[H]
\centering
\caption{Design points within each base system (16 per base, 80 total).}
\label{tab:designpoints}
\begin{tabular}{@{}ccccc@{}}
\toprule
\textbf{DP} & $\beta$ & $R_v$ &
\texttt{Branch\_Dia} & \texttt{branch\_velocity} \\
\midrule
1  & 0.25 & 0.5 & 115\,mm &  5\,m/s \\
2  & 0.25 & 1.0 & 115\,mm & 10\,m/s \\
3  & 0.25 & 2.0 & 115\,mm & 20\,m/s \\
4  & 0.25 & 4.0 & 115\,mm & 40\,m/s \\
5  & 0.50 & 0.5 & 230\,mm &  5\,m/s \\
6  & 0.50 & 1.0 & 230\,mm & 10\,m/s \\
7  & 0.50 & 2.0 & 230\,mm & 20\,m/s \\
8  & 0.50 & 4.0 & 230\,mm & 40\,m/s \\
9  & 0.75 & 0.5 & 345\,mm &  5\,m/s \\
10 & 0.75 & 1.0 & 345\,mm & 10\,m/s \\
11 & 0.75 & 2.0 & 345\,mm & 20\,m/s \\
12 & 0.75 & 4.0 & 345\,mm & 40\,m/s \\
13 & 1.00 & 0.5 & 460\,mm &  5\,m/s \\
14 & 1.00 & 1.0 & 460\,mm & 10\,m/s \\
15 & 1.00 & 2.0 & 460\,mm & 20\,m/s \\
16 & 1.00 & 4.0 & 460\,mm & 40\,m/s \\
\bottomrule
\end{tabular}
\end{table}


\section{Update Configuration}

For each base system:

\begin{enumerate}
  \item Right-click the \textbf{Parameter Set} bar $\rightarrow$
        \textbf{Properties}.
  \item Set \textbf{Design Point Update} to \textbf{Update Sequentially}
        (most stable for transient DDES).
  \item Set \textbf{Update Action} to \textbf{Update Full Project}
        (ensures geometry, mesh, and solver all update for each design
        point).
  \item Click \textbf{Update All Design Points} to begin the automated
        sweep for this base.
  \item After the 16 design points complete, move to the next base
        system and repeat.
\end{enumerate}

\medskip
\noindent\fbox{\parbox{\dimexpr\textwidth-2\fboxsep-2\fboxrule}{%
\textbf{Computation time estimate:} Each transient DDES simulation
of 25\,000 time steps with 20 iterations each requires $\sim$500\,000
solver iterations.  On an 8-core workstation, expect $\sim$16--32
hours per design point (depending on mesh size).  Each base (16
points) requires 256--512 core-hours.  The full study across all
five bases (80 design points) requires 1280--2560 core-hours.

\textbf{Parallelisation tip:} Because the five base systems are
independent Workbench systems, they can run on separate machines
simultaneously, reducing wall-clock time by up to 5$\times$.}}


% ════════════════════════════════════════════════════════════════
%  CHAPTER 10 — Post-Processing and Validation
% ════════════════════════════════════════════════════════════════
\chapter{Post-Processing and Validation}
\label{ch:post}

\section{Hydrogen Volume Fraction (Custom Field Function)}

The reference paper~\cite{Eames2022} reports hydrogen concentration as
\textbf{volume fraction} $\phi_{\mathrm{H_2}}$, while Fluent solves for
\textbf{mass fraction} $Y_{\mathrm{H_2}}$.  To enable direct comparison,
define a custom field function for the conversion:

\begin{equation}
  \phi_{\mathrm{H_2}}
  = \frac{Y_{\mathrm{H_2}} / M_{\mathrm{H_2}}}
         {Y_{\mathrm{H_2}} / M_{\mathrm{H_2}}
          + (1 - Y_{\mathrm{H_2}}) / M_{\mathrm{CH_4}}}
  \label{eq:volfrac}
\end{equation}

\noindent In Fluent, navigate to \textbf{Custom Field Functions}
$\rightarrow$ \textbf{Define} and enter:

\begin{lstlisting}[language={},numbers=none]
(Y_H2/2.016) / (Y_H2/2.016 + (1-Y_H2)/16.043)
\end{lstlisting}

\noindent Name it \texttt{vol\_frac\_H2}.  Use this function for all
contour plots, line plots, and cross-sectional averages when comparing
against the reference data.


\section{Key Output Metrics}

For each design point, extract and record the following quantities from
the time-averaged fields:

\subsection{Coefficient of Variation (COV)}

The COV quantifies mixing uniformity on a cross-sectional
plane~\cite{Kukukova2009}:

\begin{equation}
  \text{COV} = \frac{\sigma_Y}{\bar{Y}}
  \label{eq:COV}
\end{equation}

where $\bar{Y}$ is the area-weighted mean H$_2$ mass fraction on the
plane and $\sigma_Y$ is the standard deviation.
A fully mixed flow has $\text{COV} = 0$.

\noindent In Fluent, compute the COV using:
\begin{enumerate}
  \item Navigate to \textbf{Results} $\rightarrow$ \textbf{Surface
        Integrals}.
  \item Select the downstream iso-surface (e.g.\ $Z = 9200$\,mm).
  \item Report both \textbf{Area-Weighted Average} and
        \textbf{Standard Deviation} of the H2 mass fraction.
  \item Compute $\text{COV} = \text{Std.\ Dev.} / \text{Mean}$.
\end{enumerate}


\subsection[Wall H2 Concentration]{Wall H\texorpdfstring{$_2$}{2} Concentration}

To assess embrittlement risk, visualise the time-averaged H$_2$ mass
fraction on the pipe wall:

\begin{enumerate}
  \item \textbf{Results} $\rightarrow$ \textbf{Contours}.
  \item \textbf{Surface:} \texttt{wall}.
  \item \textbf{Field Variable:} Species $\rightarrow$ Mean Mass
        Fraction of h2 (from time statistics).
  \item \textbf{Colormap range:} 0 to 0.05 (to highlight regions
        exceeding 5\% mass fraction).
\end{enumerate}

\noindent Record the \textbf{maximum} H$_2$ mass fraction on the wall
and its axial location downstream of the junction.  Eames et
al.~\cite{Eames2022} reported wall concentrations exceeding 40\% by
volume for top-side injection at $\beta = 0.25$---an important
validation benchmark.


\subsection{Velocity Vectors at Junction}

\begin{enumerate}
  \item Create a \textbf{Plane} at $Z = 4600$\,mm (through the junction
        centre).
  \item \textbf{Results} $\rightarrow$ \textbf{Vectors}.
  \item \textbf{Surface:} The plane created above.
  \item Colour by \textbf{H2 mass fraction} to visualise both the
        flow pattern and species distribution simultaneously.
\end{enumerate}


\section[y+ Verification]{\texorpdfstring{$y^+$}{y+} Verification}

After the first simulation completes:

\begin{enumerate}
  \item \textbf{Results} $\rightarrow$ \textbf{Contours}.
  \item \textbf{Surface:} \texttt{wall}.
  \item \textbf{Field Variable:} Wall Fluxes $\rightarrow$
        Wall Yplus (or $y^+$).
\end{enumerate}

\noindent Verify that $y^+ \leq 1$ over the majority of the wall
surface.  Localised regions near the junction stagnation point may
have slightly elevated values ($y^+ \sim 2$--3), which is acceptable.
If $y^+$ consistently exceeds 5, the first-layer height must be reduced
and the mesh regenerated.


\section{Mesh Independence Study}
\label{sec:meshindep}

A mesh independence study must be performed for the baseline design
point ($\beta = 0.25$, $R_v = 1.0$) in \textbf{at least one} base
system (recommended: \texttt{Theta\_090}) before committing to the
full parametric sweep.  If the mesh topology differs between bases
(e.g.\ tetrahedra for \texttt{Theta\_030} vs.\ MultiZone for
\texttt{Theta\_090}), perform a separate GCI study for each distinct
mesh method.

Follow the Grid Convergence Index (GCI) procedure of
Roache~\cite{Roache1994} and Celik et al.~\cite{Celik2008}:

\subsection{Three-Mesh Protocol}

\begin{enumerate}
  \item Generate three systematically refined meshes by scaling the
        global element size and BOI element size:

  \begin{table}[H]
  \centering
  \caption{Mesh refinement levels for GCI study.}
  \label{tab:meshlevels}
  \begin{tabular}{@{}lccc@{}}
  \toprule
  \textbf{Mesh} & \textbf{Global Size} & \textbf{BOI Size} &
  \textbf{Refinement Ratio ($r$)} \\
  \midrule
  Coarse & 22.5\,mm & 7.5\,mm & --- \\
  Medium (baseline) & 15\,mm & 5\,mm & 1.5 \\
  Fine   & 10\,mm & 3.3\,mm & 1.5 \\
  \bottomrule
  \end{tabular}
  \end{table}

  \item Run each mesh through the complete transient DDES workflow
        (Hybrid Initialisation $\rightarrow$ transient run).

  \item Extract the same output quantities (COV at outlet, maximum
        wall H$_2$ concentration, pressure drop) from each mesh.
\end{enumerate}

\subsection{GCI Calculation}

For a monitored quantity $\phi$ from the three meshes
($\phi_1$ fine, $\phi_2$ medium, $\phi_3$ coarse):

\begin{enumerate}
  \item Compute the observed order of convergence:
  \begin{equation}
    p = \frac{\ln\!\left|\frac{\phi_3 - \phi_2}{\phi_2 - \phi_1}\right|}
             {\ln(r)}
    \label{eq:orderp}
  \end{equation}

  \item Compute the fine-grid GCI:
  \begin{equation}
    \text{GCI}_{\text{fine}} =
    \frac{F_s\;\left|\frac{\phi_2 - \phi_1}{\phi_1}\right|}{r^p - 1}
    \label{eq:GCI}
  \end{equation}
  where $F_s = 1.25$ is the safety factor for three or more
  grids~\cite{Roache1994}.

  \item A GCI $< 5$\,\% indicates that the medium mesh is
        sufficiently resolved for the parametric sweep.
\end{enumerate}


\section{Validation Against Reference Data}

For the baseline case ($\beta = 0.25$, $R_v$ corresponding to
$\mathcal{D} = 0.2$) in the \texttt{Theta\_090} base, compare results
against Eames et al.~\cite{Eames2022}:

\begin{itemize}
  \item \textbf{Qualitative:} H$_2$ concentration contours on the
        pipe centreline plane should show the characteristic
        stratified plume for top-side injection and improved mixing
        for alternative injection orientations.
  \item \textbf{Quantitative:} The area-weighted H$_2$ mass fraction
        at 5\,$d_1$ and 10\,$d_1$ downstream should be within 10\,\%
        of the reference values.
  \item \textbf{Maximum wall concentration:} Should reproduce the
        $>$40\,\% wall concentration reported for top-side injection
        at low dilution ratios.
\end{itemize}

\noindent Any systematic discrepancy may arise from differences between
the OpenFOAM solver (used by Eames et al.) and ANSYS Fluent, particularly
in the DES shielding function implementation.  Such differences should
be documented and discussed.

\medskip
\noindent\fbox{\parbox{\dimexpr\textwidth-2\fboxsep-2\fboxrule}{%
\textbf{Cross-base comparison:} After completing all five bases,
combine the results into a single analysis matrix.  Plot COV and maximum
wall H$_2$ concentration as functions of $\theta$, $\beta$, and $R_v$
to identify trends and interactions.  This three-dimensional parameter
space ($5 \times 4 \times 4 = 80$ points) provides a comprehensive
dataset for training surrogate models or constructing response
surfaces.}}


% ════════════════════════════════════════════════════════════════
%  BIBLIOGRAPHY
% ════════════════════════════════════════════════════════════════
\bibliographystyle{unsrtnat}
\bibliography{references}


\end{document}
